\documentclass{exam}

\newif\ifA
\newif\ifkey
\newif\ifprintselected   % required by exam_config's \ansbox

%%%%%%%%%%%%%%%%%%%%%%%%%%%%%%%%%%%%%%%%%%%%%%%%%%%%%%
%%% SET THESE IF VARIABLES
%%%%%%%%%%%%%%%%%%%%%%%%%%%%%%%%%%%%%%%%%%%%%%%%%%%%%%
% Lab name -- labs are named by the lesson they belong to, never by a
% sequential lab number (Neil, 2026-10-12): this is the L13 Lab.
\def\labLesson{13}

% is this the key or not?
\ifdefined\iskey
	\ifnum\iskey=1
		\keytrue
	\else
		\keyfalse
	\fi
\else
	%% Pick one manually
	\keyfalse % if this is NOT the key
%	\keytrue     % if this IS the key
\fi

% set up the header and footer - change for every term
\ifdefined\thisTermIn
	\def\thisterm{\thisTermIn}
\else
	\def\thisterm{Spring 2027}
\fi
%
\def\thisclass{ECE 448}

%%%%%%%%%%%%%%%%%%%%%%%%%%%%%%%%%%%%%%%%%%%%%%%%%%%%%%
%%% END SET THESE IF VARIABLES
%%%%%%%%%%%%%%%%%%%%%%%%%%%%%%%%%%%%%%%%%%%%%%%%%%%%%%
% packages and shortcuts
\input{myPackages_gr}
\usepackage[hidelinks, linkcolor=red]{hyperref}
\input{myShortcuts}
\setlength{\parindent}{0pt}
\setlength{\parskip}{6pt}

%% exam class customization
\input{exam_config}
% ECE 448 document fonts.
%
% Body text is Barlow (SIL OFL; the faces and the license are in
% latex/fonts/Barlow/). Barlow is not in TeX Live, so it is loaded as an
% OpenType face through fontspec -- which means the document must be built
% with lualatex, not pdflatex. build_practice.sh and build_lab.sh both do.
%
% Compiling with pdflatex anyway is not an error: the guard below falls back
% to the LaTeX default so the source still compiles under any engine and can
% be syntax-checked. The resulting PDF just is not the one we ship.
%
% Paths are relative to latex/, which is where both build scripts run
% (each does cd "$(dirname "$0")"). Loading by path rather than by installed
% family name means no system font install and no fontconfig, so a fresh
% clone builds a byte-comparable PDF anywhere.
%
% MATH IS DELIBERATELY LEFT IN COMPUTER MODERN. Barlow has no math companion,
% and CM math under sans body text is already the house pairing: the decks set
% words in Source Sans Pro and let MathJax render symbols in its CM-derived
% face. Moving math to a sans as well would mean unicode-math plus a sans math
% font (Fira Math being the plausible one) and re-checking every equation in
% the course for symbol coverage. That is a separate decision, not a side
% effect of this one.

\usepackage{iftex}
\ifLuaTeX
	\usepackage{fontspec}
	\defaultfontfeatures{Ligatures=TeX}
	\setmainfont{Barlow}[
		Path           = fonts/Barlow/,
		Extension      = .ttf,
		UprightFont    = *-Regular,
		ItalicFont     = *-Italic,
		BoldFont       = *-Bold,
		BoldItalicFont = *-BoldItalic,
	]
	\setsansfont{Barlow}[
		Path           = fonts/Barlow/,
		Extension      = .ttf,
		UprightFont    = *-Regular,
		ItalicFont     = *-Italic,
		BoldFont       = *-SemiBold,
		BoldItalicFont = *-BoldItalic,
	]
	% microtype earns its place here specifically because of Barlow. Its
	% glyphs are narrower than Computer Modern's, so justified prose sets
	% tighter and TeX has to squeeze interword space to fit; two paragraphs
	% in the L02/L03 sets went overfull on the first Barlow build. Protrusion
	% and font expansion absorb that -- both are lualatex features, so this
	% belongs inside the guard with everything else.
	\usepackage{microtype}

	% Mono stays Latin Modern: Barlow ships no monospaced face, and verbatim
	% and code need one that actually is fixed-pitch.
\fi
   % Barlow body text (needs lualatex)
\printselectedfalse      % full worked solutions in the key, not answers-only
\CorrectChoiceEmphasis{\color{red}}
\SolutionEmphasis{\color{red}}
\renewcommand\questionlabel{\textbf{\thequestion.}}
\renewcommand{\thepartno}{\alph{partno}}
\renewcommand{\partlabel}{(\thepartno)}

% Fillable table cell: blank in the student copy, the expected value in red on
% the KEY. Fixed-width centered columns keep the blank cells writable.
\newcolumntype{K}[1]{>{\centering\arraybackslash}p{#1}}
\newcommand{\kv}[1]{\ifkey{\color{red}#1}\fi\rule[-0.10in]{0pt}{0.34in}}

% Answers are not in this public file. Every solution, \kv value, and \ansbox
% value is \keysol{<id>}, defined in the private ECE448-faculty repo at
% lab-keys/src/ECE448_Lab_L13_DTMF_solutions.tex; build_lab.sh puts that
% folder on TEXINPUTS (KEYDIR) for the KEY. A KEY build without it stops.
\long\def\defkeysol#1#2{\expandafter\long\expandafter\def\csname keysol@#1\endcsname{#2}}
\def\keysol#1{\ifcsname keysol@#1\endcsname\csname keysol@#1\endcsname\fi}
\ifkey
	\IfFileExists{ECE448_Lab_L13_DTMF_solutions.tex}%
		{\input{ECE448_Lab_L13_DTMF_solutions}}%
		{\errmessage{KEY build needs ECE448_Lab_L13_DTMF_solutions.tex from ECE448-faculty/lab-keys/src (set KEYDIR)}}
\fi

% print the answers if this is the key
\ifkey
	\printanswers % if this IS the key
\else
	\noprintanswers  %if this is NOT the key
\fi

% print the header and footer
% need to print key on top if this is the key
\ifkey
	\firstpageheader{}{\color{red}{ KEY}}{}
	\runningheader{\thisclass\hspace{1pt} -- \thisterm }{L\labLesson\ Lab\hspace{1pt}\color{red}{ KEY}}{\thepage\hspace{1pt} of\hspace{1pt} \numpages}
	\firstpagefooter{}{}{}
	\runningfooter{}{}{}
\else
	\firstpageheader{}{}{}
	\runningheader{\thisclass \hspace{1pt} -- \thisterm }{L\labLesson\ Lab\hspace{1pt}}{\thepage\hspace{1pt} of\hspace{1pt} \numpages}
	\firstpagefooter{}{}{}
	\runningfooter{}{}{}
\fi

\runningheadrule




\begin{document}
\pagestyle{headandfoot}   % exam class

%%%%%%%%%%%%%%%%%%%%%%%%%%%%%%%%%%%%%%%%%%%%%%%%%%%%%%
% title block
%%%%%%%%%%%%%%%%%%%%%%%%%%%%%%%%%%%%%%%%%%%%%%%%%%%%%%

\begin{center}
USAF ACADEMY \\
DEPARTMENT OF ELECTRICAL AND COMPUTER ENGINEERING \\
LESSON \labLesson\ LAB - DTMF \\
\textbf{Lab Packet} \\
\thisterm \\[4mm]
\end{center}

\vspace{-4pt}
\textbf{Name:} \rule[-2pt]{2.6in}{0.4pt} \hfill
\textbf{Section:} \rule[-2pt]{0.9in}{0.4pt} \hfill
\textbf{Date:} \rule[-2pt]{1.1in}{0.4pt}

\vspace{4pt}

\fullwidth{\textbf{Learning Objective 6.2: } I can decode DTMF tones from a captured audio
signal.}

\vspace{2pt}

This packet is the L\labLesson\ Lab, in three parts, and you turn it in once, complete,
before Lesson 14. Do \textbf{Part 1} on your own before the lab period: run the receiver
and the link, and work out the block settings you will use at the bench. \textbf{Part 2} is your own
measurements, recorded in class on your own laptop. Do \textbf{Part 3} on your own after
the lab: explain what you saw, and answer the analysis questions.

\textbf{Grading.} One submission, graded against Objective 6.2 on the course's
1 / 0 scale, with revision allowed under course policy. Paste screenshots where a
question asks for one.

\textbf{Transmitting.} Nothing in this lab transmits. Question 7 uses your phone's
speaker and your laptop's microphone.

The flowgraphs are on the L13 lesson page, \texttt{L13-dtmf-rx.grc},
\texttt{L13-dtmf-link.grc}, and \texttt{L13-dtmf-keypad.grc}, and on the L13 Lab page,
\texttt{L13-dtmf-mic.grc}, with the captures \texttt{L13-capture.wav} and
\texttt{L13-lab-capture.wav}. Both captures are synthetic: \texttt{L13-dtmf-capture.grc}
dialed them and added noise. Every number below assumes the flowgraphs' settings:
$f_s = 8$ kHz, Goertzel blocks of $N = 205$ samples unless a question changes $N$, and the
decoder's rules from the lesson: the best row and column bins hold 35\% of a block's
energy, each beats its group by 6 dB, the twist is within $-10$ to $+6$ dB, and three
blocks in a row agree.

\begin{center}
\renewcommand{\arraystretch}{1.2}
\begin{tabular}{| c | c | c | c | c |}
\hline
 & 1209 Hz & 1336 Hz & 1477 Hz & 1633 Hz \\\hline
697 Hz & 1 & 2 & 3 & A \\\hline
770 Hz & 4 & 5 & 6 & B \\\hline
852 Hz & 7 & 8 & 9 & C \\\hline
941 Hz & * & 0 & \# & D \\\hline
\end{tabular}
\end{center}

\begin{questions}

\fullwidth{\large\textbf{Part 1: Pre-Lab} \normalsize(individual, before the lab period)}

%%%%%%%%%%%%%%%%%%%%%%%%%%%%%%%%%%%%%%%%%%%%%%%%%%%
\question \textbf{Run the receiver.} Put every L13 file in one folder. Open
\texttt{L13-dtmf-rx.grc} in GRC and run it on \texttt{L13-capture.wav}, the default.

\begin{parts}

	\begin{minipage}[t]{0.58\linewidth}
		\part After about 10 s, record the Digits and Numbers labels.
		\begin{solution}[0.5in]\keysol{sol1}\end{solution}
	\end{minipage}\hfill%
	\begin{minipage}[t]{0.37\linewidth}
		\raggedleft \vspace{12pt}
		\ansbox[1.6in]{\keysol{ans1}}
	\end{minipage}

	\part Paste a screenshot of the running flowgraph's window.
		\begin{solution}[1.6in]\keysol{sol2}\end{solution}

\end{parts}

%%%%%%%%%%%%%%%%%%%%%%%%%%%%%%%%%%%%%%%%%%%%%%%%%%%
\newpage
\question \textbf{Settings and a first look at noise.}

\begin{parts}

	\part Choose a number of at least seven digits that ends in \#, such as a voicemail
	code. Write each key's row and column frequency from the grid.
		\begin{solution}[0.8in]\keysol{sol3}\end{solution}

	\begin{minipage}[t]{0.58\linewidth}
		\part The bench changes the Goertzel Length $N$. Find the block time $N/f_s$ and the
		bin spacing $f_s/N$, which is also the distance to the first null, for $N = 102$ and $N = 410$. For each, say whether a 50 ms tone
		can fill a whole block.
		\begin{solution}[1.2in]\keysol{sol4}\end{solution}
	\end{minipage}\hfill%
	\begin{minipage}[t]{0.37\linewidth}
		\raggedleft \vspace{12pt}
		102: \ansbox[1.3in]{\keysol{ans2}}

		\vspace{4pt}
		410: \ansbox[1.3in]{\keysol{ans3}}
	\end{minipage}

	\begin{minipage}[t]{0.58\linewidth}
		\part Open and run \texttt{L13-dtmf-link.grc}. Set the Twist slider to $-6$ dB, so
		the column tone arrives weaker. Then lower the SNR slider from 20 dB, 1 dB at a time,
		waiting about 15 s at each step. Record the SNR at which the Numbers label first shows
		a number missed or wrong, and describe what the spectrum looks like there.
		\begin{solution}[1.2in]\keysol{sol5}\end{solution}
	\end{minipage}\hfill%
	\begin{minipage}[t]{0.37\linewidth}
		\raggedleft \vspace{12pt}
		SNR \ansbox[1.1in]{\keysol{ans4}}
	\end{minipage}

\end{parts}

%%%%%%%%%%%%%%%%%%%%%%%%%%%%%%%%%%%%%%%%%%%%%%%%%%%
\newpage
\fullwidth{\large\textbf{Part 2: Measurements} \normalsize(in class)}

\question \textbf{Build the keypad and its bins.} Start GRC with an empty canvas. Do not open
\texttt{L13-dtmf-keypad.grc}: build the keypad yourself as the lab page lists, then add the
receiver's eight Goertzel bins.

\begin{parts}

	\part Paste a screenshot of your flowgraph on the GRC canvas and of its spectrum with
	one key held. Every connection must be black, not red.
		\begin{solution}[1.6in]\keysol{sol6}\end{solution}

	\part For three keys, read both spectrum peaks and compare them with the grid.

	\begin{center}
	\renewcommand{\arraystretch}{1.25}
	\begin{tabular}{| l | K{1.1in} | K{1.1in} | K{1.1in} |}
	\hline
	Key & & & \\\hline
	Row peak (Hz) & & & \\\hline
	Column peak (Hz) & & & \\\hline
	\end{tabular}
	\end{center}

		\begin{solution}[0.5in]\keysol{sol7}\end{solution}

	\part Add the eight Goertzel bins and the DTMF Decoder to your keypad, as the lab page
	lists. Hold one key and paste a screenshot of both time sinks. Which traces rise, and how
	high?
		\begin{solution}[1.4in]\keysol{sol8}\end{solution}

	\part Press each key of your Question 2(a) number. Record the keys the decoder reports,
	and say whether they match.
		\begin{solution}[0.7in]\keysol{sol9}\end{solution}

\end{parts}

%%%%%%%%%%%%%%%%%%%%%%%%%%%%%%%%%%%%%%%%%%%%%%%%%%%
\newpage
\question \textbf{Decode the lab capture.} Run \texttt{L13-dtmf-rx.grc} on
\texttt{L13-lab-capture.wav}, as the lab page lists.

\begin{parts}

	\begin{minipage}[t]{0.58\linewidth}
		\part Record the number and paste a screenshot.
		\begin{solution}[1.2in]\keysol{sol10}\end{solution}
	\end{minipage}\hfill%
	\begin{minipage}[t]{0.37\linewidth}
		\raggedleft \vspace{12pt}
		\ansbox[1.6in]{\keysol{ans5}}
	\end{minipage}

	\begin{minipage}[t]{0.58\linewidth}
		\part Read the row and column peak heights. Which group arrives weaker, and about
		what fraction of the other group's height does it reach?
		\begin{solution}[1.0in]\keysol{sol11}\end{solution}
	\end{minipage}\hfill%
	\begin{minipage}[t]{0.37\linewidth}
		\raggedleft \vspace{12pt}
		row \ansbox[1.0in]{\keysol{ans6}}

		\vspace{4pt}
		column \ansbox[1.0in]{\keysol{ans7}}

		\vspace{4pt}
		weaker \ansbox[1.0in]{\keysol{ans8}}
	\end{minipage}

	\begin{minipage}[t]{0.58\linewidth}
		\part Find the time from one dialing to the next and estimate the time per key,
		tone plus gap.
		\begin{solution}[1.0in]\keysol{sol12}\end{solution}
	\end{minipage}\hfill%
	\begin{minipage}[t]{0.37\linewidth}
		\raggedleft \vspace{12pt}
		period \ansbox[1.0in]{\keysol{ans9}}

		\vspace{4pt}
		per key \ansbox[1.0in]{\keysol{ans10}}
	\end{minipage}

\end{parts}

%%%%%%%%%%%%%%%%%%%%%%%%%%%%%%%%%%%%%%%%%%%%%%%%%%%
\newpage
\question \textbf{The decoder's limits.} Run \texttt{L13-dtmf-link.grc} as the lab page
lists. At each setting, wait about 30 s and record the Numbers label as correct of
decoded. A count of 0 of 0 means that no number decoded.

\begin{parts}

	\part Twist, with all 16 keys as the number:

	\begin{center}
	\renewcommand{\arraystretch}{1.25}
	\begin{tabular}{| l | K{0.42in} | K{0.42in} | K{0.42in} | K{0.42in} | K{0.42in} | K{0.42in} | K{0.42in} | K{0.42in} |}
	\hline
	Twist (dB) & $-11$ & $-10$ & $-9$ & $-8$ & $+4$ & $+5$ & $+6$ & $+7$ \\\hline
	Correct & \kv{\keysol{kv1}} & \kv{\keysol{kv2}} & \kv{\keysol{kv3}} & \kv{\keysol{kv4}} & \kv{\keysol{kv5}} & \kv{\keysol{kv6}} & \kv{\keysol{kv7}} & \kv{\keysol{kv8}} \\\hline
	\end{tabular}
	\end{center}

		\begin{solution}[0.4in]\keysol{sol13}\end{solution}

	\part Tone length, with 60 ms gaps:

	\begin{center}
	\renewcommand{\arraystretch}{1.25}
	\begin{tabular}{| l | K{0.62in} | K{0.62in} | K{0.62in} | K{0.62in} | K{0.62in} | K{0.62in} |}
	\hline
	Tone (ms) & 25 & 28 & 30 & 33 & 35 & 40 \\\hline
	Correct & \kv{\keysol{kv9}} & \kv{\keysol{kv10}} & \kv{\keysol{kv11}} & \kv{\keysol{kv12}} & \kv{\keysol{kv13}} & \kv{\keysol{kv14}} \\\hline
	\end{tabular}
	\end{center}

		\begin{solution}[0.4in]\keysol{sol14}\end{solution}

	\part Gap length, with 60 ms tones:

	\begin{center}
	\renewcommand{\arraystretch}{1.25}
	\begin{tabular}{| l | K{0.62in} | K{0.62in} | K{0.62in} | K{0.62in} | K{0.62in} | K{0.62in} |}
	\hline
	Gap (ms) & 10 & 15 & 20 & 25 & 30 & 40 \\\hline
	Correct & \kv{\keysol{kv15}} & \kv{\keysol{kv16}} & \kv{\keysol{kv17}} & \kv{\keysol{kv18}} & \kv{\keysol{kv19}} & \kv{\keysol{kv20}} \\\hline
	\end{tabular}
	\end{center}

		\begin{solution}[0.6in]\keysol{sol15}\end{solution}

\end{parts}

%%%%%%%%%%%%%%%%%%%%%%%%%%%%%%%%%%%%%%%%%%%%%%%%%%%
\newpage
\question \textbf{Block length.} Change the Variable \texttt{N} in
\texttt{L13-dtmf-link.grc} as the lab page lists, and record the Numbers label.

\begin{center}
\renewcommand{\arraystretch}{1.25}
\begin{tabular}{| l | l | K{1.3in} |}
\hline
$N$ & Setting & Correct \\\hline
102 & 22 ms tones, 60 ms gaps & \kv{\keysol{kv21}} \\\hline
102 & 25 ms tones, 60 ms gaps & \kv{\keysol{kv22}} \\\hline
102 & 30 ms tones, 60 ms gaps & \kv{\keysol{kv23}} \\\hline
102 & 50 ms tones, $+1.5$\% error & \kv{\keysol{kv24}} \\\hline
102 & 50 ms tones, $+2.5$\% error & \kv{\keysol{kv25}} \\\hline
102 & 50 ms tones, $+3.5$\% error & \kv{\keysol{kv26}} \\\hline
410 & 50 ms tones and gaps & \kv{\keysol{kv27}} \\\hline
410 & 100 ms tones, $+0.5$\% error & \kv{\keysol{kv28}} \\\hline
410 & 100 ms tones, $+1.0$\% error & \kv{\keysol{kv29}} \\\hline
\end{tabular}
\end{center}

	\begin{solution}[0.8in]\keysol{sol16}\end{solution}

%%%%%%%%%%%%%%%%%%%%%%%%%%%%%%%%%%%%%%%%%%%%%%%%%%%
\newpage
\question \textbf{A real phone.} Run \texttt{L13-dtmf-mic.grc} as the lab page lists.

\begin{parts}

	\part Record the number you pressed and the Digits label, and paste a screenshot of the
	window while you press one key.
		\begin{solution}[1.5in]\keysol{sol17}\end{solution}

	\part Press keys as fast as you can. What does the decoder miss?
		\begin{solution}[0.9in]\keysol{sol18}\end{solution}

	\part Talk, whistle, or play music near the microphone. Record any false digits.
		\begin{solution}[0.9in]\keysol{sol19}\end{solution}

\end{parts}

%%%%%%%%%%%%%%%%%%%%%%%%%%%%%%%%%%%%%%%%%%%%%%%%%%%
\newpage
\fullwidth{\large\textbf{Part 3: Writeup} \normalsize(individual, after the lab)}

\question \textbf{Analysis.}

\begin{parts}

	\part Lesson 13 says DTMF tones last 40 ms or longer and gaps 20 ms or longer, and the
	decoder's twist test allows $-10$ to $+6$ dB. Compare each with what you saw in
	Question 5. Where does the decoder start to miss digits, and why there?
		\begin{solution}[1.5in]\keysol{sol20}\end{solution}

	\part Use Question 6 to explain what a shorter and a longer Goertzel block gain and
	lose, and why $N = 205$ is the usual choice at 8 kHz.
		\begin{solution}[1.5in]\keysol{sol21}\end{solution}

	\part Compare the phone with the synthetic captures, and propose one change to the
	decoder for a noisy room.
		\begin{solution}[1.5in]\keysol{sol22}\end{solution}

\end{parts}

\end{questions}

\vspace{1.0em}
\noindent\textbf{Documentation:} \rule[-2pt]{0.6\linewidth}{0.4pt}

\end{document}
